%! Modeling with Quadratics — Unit 2, Lesson 2.7 — Lesson Plan
\documentclass[10pt]{article}
\usepackage{algebra2-article}
\usepackage{algebra2-boxes}
\usepackage{graphicx}   % -article does NOT load graphicx; needed for thumbnails/screenshots
\graphicspath{{images/}}
\newcommand{\TallMath}[1]{$\displaystyle #1\rule[-1.4em]{0pt}{3.2em}$}
\newcommand{\CourseName}{Algebra 2: Shepherd}
\newcommand{\SchoolYear}{2026--2027}
\newcommand{\MeetingLength}{55 minutes}
\newcommand{\UnitNumberName}{Unit 2: Quadratic Functions \quad}
\newcommand{\LessonNumberName}{Lesson 2.7: Modeling with Quadratics}

\begin{document}
\begin{center}
    {\Large\bfseries \CourseName: \SchoolYear \\
    {\small\bfseries \UnitNumberName \LessonNumberName}}
\end{center}

\begin{tcolorbox}[colback=forestbg, colframe=forest, boxrule=0.9pt, arc=2mm,
    left=2mm, right=2mm, top=1.5mm, bottom=1.5mm]
    \textbf{Primary Objective:} In the unit's capstone lesson, students stop being handed a quadratic and instead
    \emph{build} one from a real situation --- a \textbf{projectile}, a \textbf{maximum-area} enclosure, or a
    \textbf{revenue-optimization} story --- and then read the answer to the real question off the model's key
    features. The organizing idea is a \textbf{feature-to-question} map: the \textbf{$y$-intercept} is the starting
    value, a \textbf{zero} (positive root) is ``when/where it reaches $0$,'' and the \textbf{vertex} is the
    \textbf{maximum or minimum} --- its input tells you \emph{when/where} the max/min happens and its output tells you
    \emph{how much}. The algebra is entirely from Lessons 2.1--2.5 (vertex by $x=-b/2a$; zeros by factoring or the
    formula); the genuinely new work is \textbf{translating} the situation into a model, choosing which feature the
    question asks for, and \textbf{interpreting} the result in context with units --- rejecting answers that don't fit
    the situation (negative time, a width outside the feasible domain).
    \\[2pt]
    \textbf{Standards (2023 VA SOL):} \textbf{A2.EI.2a} (create a quadratic equation/function in one variable to
    model a contextual situation), \textbf{A2.EI.2d} (verify solutions and interpret them --- including the vertex as
    max/min --- in the context of the problem), and \textbf{A2.F.2d} (identify and interpret the maximum/minimum
    value of a quadratic from its vertex). Builds directly on graphing and vertex-finding (\textbf{A2.F.2a/d}, Lesson
    2.1) and solving a single quadratic by factoring / the quadratic formula (\textbf{A2.EI.2b}, Lessons 2.2 \& 2.5).
\end{tcolorbox}

\renewcommand{\arraystretch}{1.4}
\begin{skillbox}[Priority Ideas \& Skills]{goldbox}
\begin{tabularx}{\linewidth}{@{}X|X@{}}
\textbf{Priority Ideas \& Skills} & \textbf{Key Understandings (the ``why'')} \\[2pt]
\hline
\vspace{-6pt}
\begin{itemize}[leftmargin=1.1em, nosep, after=\vspace{2pt}]
  \item Translate a projectile / area / revenue situation into a quadratic model.
  \item Read a model's \textbf{$y$-intercept} as the starting value and its \textbf{zero} as ``reaches $0$.''
  \item Find the \textbf{vertex} ($x=-b/2a$, then evaluate) and name it as the \textbf{max or min}.
  \item Decide \emph{which} feature answers the question (starting value / when-it-lands / the best value).
  \item Interpret every answer \textbf{in context with units}; reject values outside the feasible domain.
  \item Solve the model with factoring or the quadratic formula --- no new equation-solving.
\end{itemize}
&
\vspace{-6pt}
\begin{itemize}[leftmargin=1.1em, nosep, after=\vspace{2pt}]
  \item A model is only useful if the variables stand for something real --- define them first.
  \item Features are not abstract: on a height model the $y$-intercept \emph{is} the launch height.
  \item ``Maximum/minimum'' always lives at the vertex --- input = when/where, output = how much.
  \item Different questions point at different features; matching them is the modeling skill.
  \item An algebraic answer is only a real answer if it makes sense in the situation.
  \item The unit gave you every tool to solve the model; today's work is setting it up and reading it.
\end{itemize}
\end{tabularx}
\end{skillbox}

\begin{skillbox}[{Vocabulary, Concepts, \& Theorems}]{sky}
\renewcommand{\arraystretch}{1.3}
\begin{tabularx}{\linewidth}{@{}l X@{}}
\textbf{Quadratic model} & A quadratic function used to describe a real situation (height, area, revenue). \\
\textbf{Projectile model} & $h(t)=-16t^2+v_0t+h_0$ (feet, seconds): $h_0=$ launch height, $v_0=$ launch speed. \\
\textbf{Maximum / minimum value} & The output at the vertex --- the greatest (or least) value the model reaches. \\
\textbf{Optimize} & Find the input that makes a quantity as large (or small) as possible --- i.e.\ the vertex. \\
\textbf{Feasible (reasonable) domain} & The inputs that make sense in context (e.g.\ $t\ge0$, a width $>0$). \\
\textbf{Interpret in context} & State what a number means for the situation, with units. \\
\end{tabularx}
\end{skillbox}

\begin{fixedskillbox}[Activate Prior Knowledge \& Spiral Review]{forestbg}
    The Warm-Up rehearses the three moves modeling demands: \textbf{(1)} \emph{finding a vertex} of a quadratic in
    standard form via $x=-b/2a$ then evaluating (Lesson 2.1) --- the engine behind every ``maximum/minimum''
    question; \textbf{(2)} \emph{solving a quadratic} $=0$ by factoring (Lessons 2.2/2.5) --- what ``when does it hit
    the ground?'' collapses to; and \textbf{(3)} \emph{evaluating and reading a height function} --- $h(0)$ as the
    start and the vertex's output as the maximum. Debrief item~3 last: ``the vertex isn't just a point --- its
    $y$-value is the \emph{most} the ball ever climbs.''
\end{fixedskillbox}

\begin{skillbox}[Hook]{forestbg}
    Project the unit's projectile: a ball thrown \emph{up} at $32$ ft/s from a $48$-ft ledge, $h(t)=-16t^2+32t+48$.
    Ask three questions students already have the tools to answer but haven't yet \emph{named}: ``How high does it
    start?'' ($h(0)=48$, the $y$-intercept), ``What's the highest it gets, and when?'' (the vertex $(1,64)$), and
    ``When does it hit the ground?'' (the positive zero, $t=3$). The pivot: \emph{every} question about the story is
    really a question about a \emph{feature} of the parabola --- and you learned to find all of them in Lessons
    2.1--2.5. Today you supply the model, then read the story off it.
\end{skillbox}

\begin{skillbox}[Lesson]{forestbg}
\begin{multicols}{2}
\textbf{1. Projectile motion.} Establish the model $h(t)=-16t^2+v_0t+h_0$ on the anchor $h(t)=-16t^2+32t+48$. Map
each question to a feature: start height $=h(0)=48$ ($y$-int); maximum height $=$ vertex, $t=-b/2a=1$, $h(1)=64$;
landing $=$ positive zero, $-16t^2+32t+48=0\Rightarrow t^2-2t-3=0\Rightarrow t=3$ (reject $t=-1$). Stress
\emph{units} and \emph{reject the impossible root}.

\textbf{2. Maximum area.} A pen against a barn uses $40$ ft of fence on \emph{three} sides. Width $x$ (twice),
length $40-2x$. Area $A(x)=x(40-2x)=-2x^2+40x$. Vertex $x=-b/2a=10\Rightarrow A=200$; dimensions $10\times20$.
The vertex is the \emph{maximum area}; the feasible domain is $0<x<20$.

\columnbreak

\textbf{3. Optimize revenue.} A theater sells $200$ tickets at \$8; each \$1 increase loses $20$ buyers. Price
$8+x$, attendance $200-20x$, revenue $R(x)=(8+x)(200-20x)=-20x^2+40x+1600$. Vertex $x=1\Rightarrow$ price \$9,
revenue \$1620. ``Best price'' is an \emph{optimization}: it lives at the vertex.

\textbf{4. The strategy (post it).} (i) Define the variable. (ii) Write the quadratic model. (iii) Decide which
feature the question wants: start $\to$ $y$-int; ``reaches $0$'' $\to$ zero; ``most/least/best'' $\to$ vertex.
(iv) Find it (Lessons 2.1--2.5). (v) \emph{Interpret with units}, rejecting values outside the feasible domain.
\end{multicols}
\end{skillbox}

\begin{skillbox}[Explicit Instruction: From a Story to the Vertex]{forestbg}
\begin{tabularx}{\linewidth}{@{}p{0.46\linewidth} X@{}}
\textbf{Strategy (post and refer back):}
\begin{enumerate}[leftmargin=1.4em, nosep, itemsep=2pt]
  \item Define the variable (with units).
  \item Write the quadratic model.
  \item Ask what feature the question wants.
  \item Find it: $y$-int $=c$; zeros by factoring/formula; vertex by $x=-b/2a$ then evaluate.
  \item Interpret in context, with units.
  \item Reject answers outside the feasible domain.
\end{enumerate}
&
\textbf{Worked example.} A rocket: $h(t)=-16t^2+64t$ (ft, s). \emph{Highest point?}
\begin{itemize}[leftmargin=1.2em, nosep, itemsep=1pt]
  \item Variable: $t=$ seconds after launch.
  \item ``Highest'' $\to$ vertex. $t=-\tfrac{64}{2(-16)}=2$.
  \item $h(2)=-16(4)+64(2)=64$.
  \item Max height $64$ ft, reached at $t=2$ s.
  \item Lands: $-16t^2+64t=0\Rightarrow t=0,4$; it returns at $t=4$ s (reject nothing --- both are real, $t=0$ is launch).
\end{itemize}
\end{tabularx}
\end{skillbox}

\begin{skillbox}[Active Monitoring]{redbox}
Circulate during the notes practice and Tier work. Watch for and correct:
\begin{itemize}[leftmargin=1.2em, nosep]
  \item reporting the vertex's \emph{input} when the question wants the \emph{output} (or vice versa) --- ``when'' vs.\ ``how high'';
  \item forgetting the length is $40-2x$ (not $40-x$) when the barn wall replaces the fourth side;
  \item keeping a \emph{negative} root as a landing time, or a width outside the feasible domain;
  \item stopping at the model and never answering the actual question (no interpretation, no units);
  \item setting the revenue model up as price\,$\times$\,\emph{fixed} sales instead of price\,$\times$\,\emph{changing} demand;
  \item confusing the $y$-intercept (start) with the vertex (max/min).
\end{itemize}
Cold-call prompts: ``Does the question want \emph{when} or \emph{how much}?'' ``Which feature is that?'' ``What are
the units of your answer?'' ``Is a width of $30$ possible with only $40$ ft of fence?'' ``Why do we throw out $t=-1$?''
\end{skillbox}

\begin{skillbox}[Group Work \& Differentiation]{redbox}
\textbf{Model It!} activity (tiered; mirrors the notes).
\begin{multicols}{3}
\textbf{Tier R --- Remediate.} Read a pre-drawn projectile graph (start / max / landing); match questions to features
(start $\to$ $y$-int, best $\to$ vertex, reaches $0$ $\to$ zero); interpret a \emph{given} area model's vertex.

\columnbreak
\textbf{Tier A --- Approaching.} Build and solve one projectile, one maximum-area (three sides), and one
maximum-revenue model; find the vertex and interpret it with units.

\columnbreak
\textbf{Tier E --- Extension.} Optimize with a feasible-domain statement; a ``same-height'' object comparison that
ties back to systems (Lesson 2.6); verify and interpret.
\end{multicols}
\end{skillbox}

\begin{skillbox}[Individual Work \& Assessment]{redbox}
\textbf{Exit Ticket} (independent, no notes): read a projectile model for its \textbf{maximum height} (vertex output)
and its \textbf{landing time} (positive zero); find a \textbf{maximum area} from a fencing model; and an
\textbf{SOL-style multiple-choice} item on which feature represents the maximum height. Collect and sort into ``got
it / found the model but read the wrong feature / no interpretation-or-units / kept an impossible value'' to decide
whether the unit review opens with a modeling re-teach.
\end{skillbox}

\begin{skillbox}[Reinforcement \& Extension]{goldbox}
\textbf{Homework:} build and interpret a projectile model (start height, max height \& time, landing); a
maximum-area (three-sides) model; a maximum-revenue model; an interpret-the-features item (what the $y$-intercept,
vertex, and zero each mean); and a ``find the error'' item on the $40-2x$ trap. \textbf{Extension:} a feasible-domain
optimization and a same-height comparison (systems tie-back). \textbf{Preview:} this closes Unit 2 --- next is the
\emph{Unit 2 Review \& Test}, pulling together graphing, factoring, the quadratic formula, complex numbers, systems,
and today's modeling.
\end{skillbox}
% ── Teacher notes (migrated from the component keys) ──
% One per component, titled for it. They live here rather than in the _key
% files so that every component is the same length blank and keyed.

\begin{teachernote}[Warm-Up]
Each item is a modeling sub-skill. Item~1 is the vertex engine, $x=-\frac{b}{2a}=-\frac{4}{-2}=2$, $f(2)=5$; because
$a=-1<0$ the parabola opens \emph{down}, so $(2,5)$ is a \emph{maximum} --- this is exactly ``how high, and when'' for
a projectile. Item~2 is what ``when does it hit the ground?'' collapses to: $-16t(t-3)=0\Rightarrow t=0,3$ (note the
GCF and the easy-to-miss $t=0$, the launch instant). Item~3 builds the reading habit: $h(0)=0$ means the ball leaves
from the ground; its vertex ($t=1.5$, $h=36$) would be the \emph{maximum} height. Debrief item~3 last, connecting
``vertex $y$-value $=$ greatest height'' to today's projectile work.
\end{teachernote}

\begin{teachernote}[Guided Notes]
The spine: \emph{every} question about a story is a question about a \emph{feature}, and students already know how to
find all of them. In \S1 hammer the units and the ``reject $t=-1$'' habit --- and the distinction between the vertex's
\emph{input} ($t=1$, \emph{when}) and \emph{output} ($64$ ft, \emph{how high}). \S2's classic error is length $=40-x$
instead of $40-2x$; walk the fence-count on the diagram. \S3 sets revenue as price\,$\times$\,\emph{changing} demand,
not fixed sales; the vertex is the ``best price.'' Practice \#1's firework lands at $t=5$ ($t^2-4t-5=0$); \#2 maximizes
at the vertex $x=15$ ($15\times30=450$); \#3's $R$ peaks at $x=15$, $R=2250$. If time is short, \#1 and \#2 are the
must-dos --- they carry ``vertex output $=$ max'' and the $40-2x$ setup.
\end{teachernote}

\begin{teachernote}[Group Activity]
Tier R stays graphical/verbal: reading start/max/landing off the curve, matching questions to features, and reading a
\emph{given} vertex --- no algebra. Tier A is the build-and-solve engine; insist on units and on the vertex output vs.\
input distinction (\#1: max is $144$ ft, \emph{at} $t=1$). \#2's frequent error is $A=x(80-x)$ --- the barn wall means
$80-2x$. \#3 sets revenue as price\,$\times$\,\emph{changing} demand ($R=-15x^2+120x+3600$, vertex $x=4$, price \$16).
Tier E \#1 adds the feasible domain $0<x<18$ (an honest ``interpret in context'' beat); \#2 reuses substitution from
Lesson 2.6 --- the two objects meet at $t=2$, both at $32$ ft. Early finishers: ask them for Object A's \emph{own}
maximum height ($t=1.5$, $36$ ft) and whether the objects are ever at the same height while A is falling.
\end{teachernote}

\begin{teachernote}[Exit Ticket]
Sort into ``got it / found the model but read the wrong feature / no interpretation-or-units / kept an impossible
value.'' Telltales: 1(a) giving $t=2$ (the \emph{when}) as the height instead of $144$ ft, or vice versa; 1(b)
keeping $t=-1$; \#2 stopping at $x=25$ without the area, or mis-setting the model; \#3 choosing \textbf{B} (input vs.\
output) or \textbf{D} (the ``it lands'' default). If \#3 or 1(a) is widely missed, open the unit review with a
60-second ``vertex input $=$ when, vertex output $=$ how much'' recap.
\end{teachernote}

\begin{teachernote}[Homework]
Common slips: \#1(b) reporting $t=1$ (the \emph{when}) as the height instead of $144$ ft, or keeping $t=-2$ in (c);
\#2 using $48-x$ instead of $48-2x$ (the headline trap); \#3 setting revenue as price\,$\times$\,\emph{fixed} $500$
members rather than the changing $500-10x$; \#4 is pure interpretation --- credit answers that tie $(0,1600)$ to ``no
increase'' and $(1,1620)$ to ``one \$1 increase, the maximum.'' \#5 is the lesson's headline misconception (three
sides $\Rightarrow$ $40-2x$); the corrected max area is $200$ ft$^2$. Extension \#1 reuses Lesson 2.6 substitution
(the drones meet at $t=2$, both $64$ ft); \#2 is an honest feasible-domain statement.
\end{teachernote}

\end{document}
